\section{Related Work}
\label{sec:related}

\paragraph{State-space models for sequence modelling.}
Structured state-space models~\citep{gu2021combining,gu2022efficiently} reformulate sequence mixing as a parameterised linear dynamical system, recovering near-linear sequence-length scaling. The S4 family~\citep{gu2022efficiently} and the diagonal variants~\citep{gupta2022diagonal,gu2022parameterization} introduced practical kernels and parameterisations. Mamba~\citep{gu2023mamba} added input-dependent selectivity: the discretisation step $\Delta_t$, the input matrix $B_t$, and the readout $C_t$ are all functions of the current input, granting the SSM a content-based memory mechanism. Mamba-2~\citep{dao2024transformers} unifies the SSM update with a structured-attention view and improves hardware utilisation. We build on Mamba-1 and treat its state $h_t$, of shape $\dinner \times \dstate$, as the natural locus of the intracellular-attention mechanism we propose.

\paragraph{Hybrid SSM--attention architectures.}
Pure-SSM models lag pure-attention transformers on tasks requiring precise in-context recall~\citep{jelassi2024repeat,arora2023zoology}. Hybrid models address this by interleaving a small number of attention layers among many SSM layers. Jamba~\citep{lieber2024jamba} popularised the design at scale; Zamba~\citep{glorioso2024zamba} explored parameter-sharing across the attention layers; \citet{park2024empirical} systematically studied layer-mixing ratios and found that even a single attention layer per block of SSM layers recovers most of the lost expressivity. We adopt a 5:1 SSM-to-attention ratio with multi-query attention~\citep{shazeer2019fast} and rotary positional embeddings~\citep{su2021roformer} for the attention layers. Our two contributions are orthogonal to the choice of mixing ratio and could be applied to any of the above hybrid designs.

\paragraph{Activation steering and representation engineering.}
A growing body of work shows that LLM behaviour is modulated by a small number of linear directions in the residual stream. \citet{zou2023representation} introduced representation-engineering as a framework for reading and writing such directions. \citet{turner2023activation} and \citet{rimsky2024steering} demonstrate that adding a difference-of-means vector --- computed from contrastive prompts --- to the residual stream at inference time can steer attributes like truthfulness, refusal, and persona. Earlier work on contrast-consistent search~\citep{burns2022discovering} and probing~\citep{alain2017understanding} likewise establishes that interpretable axes are present in unmodified models. All of this work, however, applies the directions \emph{post-hoc} at inference time, with weights chosen by the user. Our hormone-routing module instead re-injects extracted directions \emph{during pretraining} through a learned router that decides per-token how to combine them; the directions themselves are kept frozen so that their geometric meaning is preserved through training.

\paragraph{Routing and mixture-of-experts.}
Mixture-of-experts (MoE) models~\citep{shazeer2017outrageously,fedus2022switch,jiang2024mixtral} use learned routers to dispatch tokens to specialised expert subnetworks. Our hormone router is structurally analogous --- a linear-then-softmax map from the hidden state to a distribution over a small, fixed set of directions --- but the ``experts'' here are not parameter banks; they are \emph{frozen direction vectors} of dimension $\dmodel$. The routing decision selects which behavioural axis to attend to, not which subnetwork to compute. This permits a richer routing signal at a vanishing parameter cost.

\paragraph{Soft prompts and prefix tuning.}
Trained prompt embeddings~\citep{lester2021power,li2021prefix} introduce a small number of trainable vectors prepended to the input; these are conceptually related to our hormones in that both use a small bank of $\dmodel$-dimensional vectors to modulate behaviour. The differences are substantive: hormone vectors are not prefix tokens, are not exposed to attention, are token-conditional rather than fixed across the sequence, and are extracted from the base model rather than initialised randomly.

\paragraph{Slot-based attention.}
``Slots'' have appeared in the visual representation-learning literature~\citep{locatello2020object,jaegle2021perceiver} as latent tokens that summarise scene content via attention. Our use of the term refers to the $\dstate$ memory cells inside an SSM state and is unrelated; the conceptual link is only that both designs apply attention over a small, learned set of summary vectors. To our knowledge, no prior work applies attention over the state dimension of a selective SSM.
